\par\medskip
\noindent 2.3.19. \underline{Otra demostración de 2.3.11}
(por Deligne).

\medskip
\noindent a) Reducción a $P=\mathbb Z/n\mathbb Z$ y a
$T=(\mathrm{Sets})$.

\medskip
\noindent b) Sea $\partial$ la flecha siguiente de $nQ$ a
$\,{}_nQ[1]$:
\[
 nQ\xleftarrow[\text{``$n$''}]{\ \sim\ }
 [\,{}_nQ\longrightarrow Q\,]
 \longrightarrow{}_nQ[1].
\]
El diagrama siguiente es anticonmutativo:
\[
\begin{tikzcd}[row sep=huge,column sep=huge]
 nQ[\sisi{o}{0}]\arrow[r,"1"]\arrow[d,"\partial"']
   &{[Q\xrightarrow{n}Q]}\\
 {}_nQ[1]\arrow[ru,"1"']&{.}
\end{tikzcd}
\]
En efecto, las flechas
\[
\begin{tikzcd}[column sep=small,row sep=large]
 {[{}_nQ}\arrow[r]\arrow[dd,"0"']&{Q]}\arrow[dd,"n"]
   &&{[{}_nQ}\arrow[r]\arrow[dd,"1"']&{Q]}\arrow[dd,"0"']
   &\text{son opuestas}
   &{[{}_nQ}\arrow[r]&{Q]}\arrow[ddl,"1"]\\
 &&\text{y}&&&\text{por medio de}\quad H\,:&&\\
 {[Q}\arrow[r,"n"]&{Q]}&&{[Q}\arrow[r,"n"]&{Q]}
   &&{[Q}\arrow[r,"n"]&{Q]}
\end{tikzcd}
\]

%%%%%%%%%%  scan idx 263  |  folio 252  %%%%%%%%%%

\medskip
\noindent c) El diagrama 2.3.10 es
\[
\begin{tikzcd}[column sep=small,row sep=small]
 &&&H^1\mathrm{RHom}
   (\mathbb Z\mathbin{\overset{\mathbb L}{\otimes}}Q,G)
   \arrow[ddd]\\
 &&&+\\
 H^1\mathrm{RHom}
   (\mathbb Z/n\mathbb Z\mathbin{\overset{\mathbb L}{\otimes}}Q,G)
   \arrow[rr]\arrow[dd]\arrow[uurrr]
 &&H^1\mathrm{RHom}
   (\mathbb Z/n\mathbb Z,\mathrm{RHom}(Q,G))
   \arrow[dd]&\\
 &&&\mathrm{Ext}^1(Q,G)\arrow[ddl]\\
 \mathrm{Hom}(\mathrm{Tor}_1(\mathbb Z/n\mathbb Z,Q),G)\arrow[d]
 &|[draw,circle,inner sep=1pt]|{-\,?}
 &\mathrm{Hom}(\mathbb Z/n\mathbb Z,\mathrm{Ext}^1(Q,G))
   \arrow[d]&\\
 \mathrm{Hom}({}_nQ,G)\arrow[rr]
 &&\mathrm{Ext}^1(nQ,G)&
\end{tikzcd}
\]
y:

\smallskip
\noindent c1) Se tiene un isomorfismo
\[
 \mathbb Z/n\mathbb Z
 \mathbin{\overset{\mathbb L}{\otimes}}Q
 \simeq[Q\xrightarrow{n}Q],
\]
y, mediante este isomorfismo, la primera flecha vertical compuesta
está inducida por
\[
 {}_nQ[1]\longrightarrow[Q\xrightarrow{n}Q].
\]

\smallskip
\noindent c2) La segunda flecha horizontal es la transpuesta de la
flecha $\partial$ mostrada arriba en b).

\smallskip
\noindent c3) Recorriendo el diagrama en sentido horario, es claro que
se obtiene la flecha inducida por
\[
 nQ[0]\longrightarrow[Q\xrightarrow{n}Q].
\]
Esto prueba 2.3.11.
